\section{One source for the PostgreSQL boundary}
\label{sec:sql-codegen}

The preceding chapter described clients generated from provider-owned OpenAPI
descriptions. The generator considered here serves a different boundary. It
derives the extension's PostgreSQL interface from BioSeqDB declarations.

A PostgreSQL extension repeats the same contract in several forms. SQL needs a
routine name, argument types, result type, volatility, and permissions. The
shared library needs a matching C entry point. C++ code must decode the same
arguments and construct the declared result. Tables, operators, casts,
configuration variables, and comments introduce similar duplication. A local
change can therefore leave individually valid artifacts that disagree at the
boundary.

BioSeqDB makes typed C++ declarations the source of these facts. Attributes
describe only information that cannot be inferred from the declaration, such
as an SQL type binding or a PostgreSQL volatility category. Function names,
parameter order, and C++ types remain ordinary language constructs. Repeated
annotations and conflicting defaults are rejected. This keeps annotations from
becoming another complete copy of the interface.

The generator performs structural compiler analysis rather than searching
source text. It observes active declarations, attributes, parameter types,
initializers, documentation, and module ownership under the configured build
profile. These records are normalized into a closed intermediate
representation. Each SQL entity receives a stable identity and explicit
dependencies. Tarjan's method classifies strongly connected components, then
a stable topological traversal emits the component graph in deterministic
order~\cite{tarjan_1972_depth_first_search}.

\begin{figure}[htbp]
  \centering
  \begin{tikzpicture}[
      font=\small,
      node distance=0.34cm,
      stage/.style={draw,rounded corners,fill=blue!7,align=center,
                    minimum height=0.72cm,text width=0.78\linewidth},
      outputstage/.style={stage,fill=green!10},
      >=Latex]
    \node[stage] (decl) {Typed C++ declarations and checked annotations};
    \node[stage,below=of decl] (ast) {Compiler-observed declarations and active build profile};
    \node[stage,below=of ast] (ir) {Normalized SQL entities, stable identities, and dependency graph};
    \node[outputstage,below=of ir] (render) {Installation SQL, C entry points,
      row codecs, configuration registration, and manifests};
    \draw[->,thick] (decl) -- (ast);
    \draw[->,thick] (ast) -- (ir);
    \draw[->,thick] (ir) -- (render);
  \end{tikzpicture}
  \caption{Source-driven generation of the PostgreSQL extension boundary.}
  \label{fig:sql-codegen-flow}
\end{figure}

\section{Generated artifacts and safety}
\label{sec:sql-codegen-artifacts}

One declaration can produce several related artifacts. Functions normally
yield installation statements, generated function-manager entry points,
conversions, documentation, and manifest records. A row declaration can also
yield tuple ordinals and codecs. An annotated configuration value yields
registration code with a typed default, range, context, and description. A
closed type registry rejects unsupported C++ and SQL combinations rather than
falling back to an untyped \texttt{Datum} convention.

Dependencies determine a stable generation order. Generated outputs are
written to the build tree and used directly as extension artifacts.

Error handling requires an additional boundary. PostgreSQL can report a
recoverable error by a non-local jump that does not run intervening C++
destructors~\cite{postgresql17_cpp_extensibility}. Conversely, a C++ exception
must not cross a PostgreSQL callback. Generated wrappers therefore translate
both directions through a common boundary. Diagnostics are copied to plain
state, C++ owners are destroyed, and only then is the PostgreSQL error raised.
Code that calls PostgreSQL below the wrapper still needs a bridge or a region
without live nontrivial C++ objects.

The model is deliberately closed. Adding a new PostgreSQL object category
requires a typed representation, validation, dependency rules, and a renderer.
Compiler-based collection also adds build cost. In return, SQL declarations,
C symbols, converters, tuple layouts, and configuration metadata are derived
from the same program rather than maintained as independent inventories.
